\documentclass[11pt,a4paper]{article}

\usepackage[utf8]{inputenc}
\usepackage[T1]{fontenc}
\usepackage[french]{babel}
\usepackage{geometry}
\geometry{top=2.5cm, bottom=2.5cm, left=2.5cm, right=2.5cm}

\usepackage{amsmath, amssymb}
\usepackage{hyperref}
\usepackage{booktabs}
\usepackage{xcolor}

% Chemins et noms de colonnes en police fixe ne se coupent pas : marge de secours
\emergencystretch=2em

\hypersetup{
    colorlinks=true,
    linkcolor=blue,
    filecolor=blue,
    urlcolor=blue
}

\title{\textbf{DOCUMENTATION TECHNIQUE ET FONCTIONNELLE}\\[0.5em]
\Large Traitement, Analyse et Prévisions du PIB Mondial}
\author{}
\date{Septembre 2026}

\begin{document}

\maketitle

\begin{abstract}
Ce document constitue la référence technique du projet de traitement des données de PIB (Gross Domestic Product). Il détaille la méthodologie d'extraction des données historiques auprès de la Banque Mondiale (2000--2024, 25 ans) et des prévisions auprès du Fonds Monétaire International (FMI WEO, jusqu'à l'horizon de l'édition en cours), l'unification des séries temporelles, le calcul des indicateurs de croissance (CAGR), ainsi que le guide du Notebook Jupyter d'analyse interactive.
\end{abstract}

\tableofcontents
\newpage

\section{Introduction et Contexte}
L'objectif de ce projet est de constituer une base de données consolidée, automatisée et complète pour l'étude du PIB mondial. Les données couvrent plus de 190 pays sur une plage continue de 31 ans:
\begin{itemize}
    \item \textbf{Données Historiques (2000 -- 2024)}: 25 ans de séries chronologiques officielles fournies par la Banque Mondiale (World Development Indicators - WDI).
    \item \textbf{Données de Prévision (2025 -- 2031)}: Projections macroéconomiques du Fonds Monétaire International (FMI - World Economic Outlook), jusqu'à l'horizon de l'édition en cours --- cinq ans après l'année de l'édition, soit 2031 pour celle d'avril 2026.
\end{itemize}

\section{Sources de Données et APIs}

\subsection{Banque Mondiale (WDI API)}
La Banque Mondiale fournit une API REST libre d'accès sans clé d'API.
\begin{itemize}
    \item \textbf{Endpoint}: \url{https://api.worldbank.org/v2/country/all/indicator/}\texttt{\{code\}}
    \item \textbf{Indicateurs extraits}:
    \begin{enumerate}
        \item \texttt{NY.GDP.MKTP.CD}: PIB Nominal en USD courants.
        \item \texttt{NY.GDP.MKTP.KD.ZG}: Taux de croissance annuel du PIB (\%).
        \item \texttt{NY.GDP.PCAP.CD}: PIB par habitant en USD courants.
        \item \texttt{NY.GDP.MKTP.PP.CD}: PIB en Parité de Pouvoir d'Achat (PPA) en USD internationaux.
        \item \texttt{NY.GDP.MKTP.KD}: PIB à prix constants de 2015, en USD.
        \item \texttt{NY.GDP.MKTP.PP.KD}: PIB à PPA, en dollars internationaux constants de 2021.
        \item \texttt{SP.POP.TOTL}: population totale.
    \end{enumerate}
\end{itemize}

\subsection{Fonds Monétaire International (IMF DataMapper API)}
Le FMI met à disposition l'API \textit{DataMapper} (World Economic Outlook - WEO).
\begin{itemize}
    \item \textbf{Endpoint}: \url{https://www.imf.org/external/datamapper/api/v1/}\texttt{\{code\}}
    \item \textbf{Indicateurs extraits}:
    \begin{enumerate}
        \item \texttt{NGDPD}: PIB nominal en Milliards USD.
        \item \texttt{NGDP\_RPCH}: Croissance réelle du PIB (\%).
        \item \texttt{PPPGDP}: PIB PPA en Milliards d'USD internationaux.
        \item \texttt{NGDPDPC}: PIB par habitant en USD.
        \item \texttt{LP}: population, en millions.
    \end{enumerate}
\end{itemize}

Pour l'évaluation des prévisions, le FMI fournit en outre les éditions passées du WEO: la \textit{WEO Historical Forecasts Database}, téléchargée à la main, et, pour les éditions les plus récentes, son API SDMX (\url{https://api.imf.org}); voir la section~\ref{sec:evaluation}.

\section{Méthodologie et Traitement des Données}

\subsection{Harmonisation des Unités}
Les données brutes de la Banque Mondiale expriment le PIB nominal en USD bruts ($10^0$), tandis que l'API du FMI exprime le PIB en Milliards d'USD ($10^9$). Le pipeline convertit l'ensemble des valeurs en \textbf{Milliards de Dollars US (Billion USD)} selon la formule:
\begin{equation}
\text{PIB}_{\text{Milliards USD}} = \frac{\text{PIB}_{\text{USD Bruts}}}{10^9}
\end{equation}

\subsection{Calcul du Taux de Croissance Annuel Composé (CAGR)}
Pour évaluer la dynamique économique de chaque pays sur la période historique (2000--2024) et la période de prévision (2024--2031), nous utilisons le CAGR (\textit{Compound Annual Growth Rate}):
\begin{equation}
\text{CAGR} = \left( \frac{V_{\text{finale}}}{V_{\text{initiale}}} \right)^{\frac{1}{n}} - 1
\end{equation}

Où $V_{\text{initiale}}$ est la valeur de départ, $V_{\text{finale}}$ la valeur d'arrivée, et $n$ le nombre d'années ($n = 24$ pour 2000--2024, $n = 7$ pour 2024--2031).

\subsection{PIB Nominal et PIB en Volume}
Le PIB nominal en dollars courants agrège trois mouvements distincts: la croissance réelle de l'activité, l'inflation domestique et la variation du taux de change face au dollar. Comparer des CAGR nominaux entre pays revient donc à comparer des grandeurs hétérogènes.

Le pipeline collecte en parallèle deux séries en volume auprès de la Banque Mondiale:
\begin{itemize}
    \item \texttt{NY.GDP.MKTP.KD}: PIB à prix constants de 2015, en dollars US.
    \item \texttt{NY.GDP.MKTP.PP.KD}: PIB à parité de pouvoir d'achat, en dollars internationaux constants de 2021.
\end{itemize}

Ces deux mesures décrivent le même indice de volume à une constante de conversion près; elles partagent donc, pour un pays donné, un taux de croissance identique.

Le FMI ne publiant pas de \emph{niveau} de PIB en volume mais uniquement un taux (\texttt{NGDP\_RPCH}), les niveaux réels de l'horizon de prévision sont reconstitués par chaînage à partir du dernier point observé:
\begin{equation}
\text{PIB}^{\text{réel}}_{y} = \text{PIB}^{\text{réel}}_{y-1} \times \left(1 + \frac{g_{y}}{100}\right)
\end{equation}
où $g_{y}$ est la croissance réelle projetée. Une année de croissance manquante interrompt le chaînage, les années suivantes restant vides plutôt qu'extrapolées.

L'écart entre CAGR nominal et CAGR réel, exprimé en points de pourcentage, isole la contribution conjointe de l'inflation et du change. Le Japon en fournit le cas limite: $-0{,}77\%$ de croissance nominale en dollars sur 2000--2024 contre $+0{,}65\%$ en volume, soit un écart de $-1{,}4$ point imputable à la dépréciation du yen. À l'inverse, la Russie affiche $9{,}3\%$ en nominal pour $3{,}1\%$ en volume.

\subsection{Recouvrement des Sources et Raccord des Niveaux}
Le WEO couvre aussi les années observées. Le pipeline conserve ces estimations du FMI en regard des valeurs de la Banque Mondiale, qui restent la référence: elles mesurent l'écart entre institutions (\texttt{Ecart\_Sources\_Pct}) sans jamais se substituer à l'observation.

Les deux institutions concordent au centième de pourcent pour la plupart des pays (écart médian de $0{,}03\%$ sur le PIB nominal de 2024) mais divergent nettement pour une quinzaine. Juxtaposer l'observation de l'une et la projection de l'autre fabriquerait alors une croissance inexistante à la jonction. Les niveaux projetés sont donc raccordés au dernier niveau observé, ce qui conserve la dynamique du FMI et le niveau de la Banque Mondiale:
\begin{equation}
\text{PIB}_{y} = \text{PIB}^{\text{BM}}_{b} \times \frac{\text{PIB}^{\text{FMI}}_{y}}{\text{PIB}^{\text{FMI}}_{b}}
\end{equation}
où $b$ est la dernière année observée. Chaque série de niveau --- nominal, PPA courante, PIB par habitant --- reçoit son propre facteur, les deux institutions n'y divergeant pas du même rapport. Appliquer à la PPA le facteur du nominal recréait la marche que le raccord doit supprimer: le Burundi y chutait de $46\%$ en 2025, quand le FMI lui projetait $+6{,}9\%$.

Les facteurs appliqués sont conservés dans les colonnes:
\begin{itemize}
    \item \texttt{Facteur\_Raccord} (nominal);
    \item \texttt{Facteur\_Raccord\_PPA} (PPA courante);
    \item \texttt{Facteur\_Raccord\_Par\_Habitant} (PIB par habitant).
\end{itemize}

\subsection{Population et PIB par Habitant}
Les modèles de trafic aérien raisonnent par habitant. La population observée vient de la Banque Mondiale (\texttt{SP.POP.TOTL}); la population projetée, du FMI (\texttt{LP}), raccordée à son dernier niveau observé par son propre facteur, comme les autres séries de niveau: les deux sources divergent de plus de $5\,\%$ pour 42 pays. Le PIB en volume par habitant (\texttt{GDP\_Real\_Per\_Capita\_USD\_2015}, et à parité \texttt{GDP\_Real\_PPP\_Per\_Capita\_Intl\_2021}) est le volume chaîné divisé par la population raccordée, sur toute la série; la synthèse en donne les niveaux et les croissances aux années de référence. La croissance par habitant peut être bien plus faible que celle du PIB: de 2024 à 2031, $2{,}0\,\%$ par an au Nigeria pour $4{,}2\,\%$ de PIB.

\subsection{Comparaison de Niveaux: la Parité de Pouvoir d'Achat}
Neutraliser le change sur la \emph{croissance} ne suffit pas à comparer des \emph{niveaux}: la série en volume reste convertie en dollars au taux de l'année de base. Or le taux de change de marché est déterminé par les flux financiers et les biens échangeables, non par ce qu'une unité de production permet d'acheter localement.

Le classement par niveau s'appuie donc en parallèle sur \texttt{NY.GDP.MKTP.PP.KD}, exprimé en dollars internationaux constants de 2021. Les colonnes \texttt{Rank\_PPA\_2024} et \texttt{Rank\_PPA\_2031} en découlent, ainsi que l'écart \texttt{Ecart\_Rang\_Nominal\_PPA\_2024}, positif lorsqu'un pays est mieux classé à parité qu'au taux de marché.

L'effet sur le classement 2024 est substantiel: la Chine devance les États-Unis, l'Inde passe du 5\textsuperscript{e} au 3\textsuperscript{e} rang, la Russie du 10\textsuperscript{e} au 4\textsuperscript{e} et l'Indonésie du 16\textsuperscript{e} au 8\textsuperscript{e}. Symétriquement, le Canada perd six rangs et le Royaume-Uni quatre. Aucune des deux bases n'est intrinsèquement préférable: le taux de marché mesure le poids financier international, la parité de pouvoir d'achat le volume de production réellement disponible.

\subsection{Distinction entre Pays et Agrégats}
Les deux sources diffusent, au même format et dans le même flux que les pays, des \textbf{agrégats} régionaux et analytiques: \textit{World}, \textit{OECD members}, \textit{Euro area}, groupes de revenu, etc. Ces entités portent elles aussi un code sur trois lettres (\texttt{WLD}, \texttt{OED}, \texttt{EUU}) et ne peuvent donc pas être écartées sur la seule longueur du code.

Le pipeline interroge à cet effet les points d'entrée de métadonnées de chaque source:
\begin{itemize}
    \item Banque Mondiale, \texttt{/v2/country}: une entité dont la région est codée \texttt{NA} est un agrégat (78 agrégats pour 217 pays).
    \item FMI, \texttt{/api/v1/countries}: tout code absent de cette liste de 241 pays est un groupe analytique (\texttt{ADVEC}, \texttt{EURO}, \texttt{MENA}\ldots).
\end{itemize}

Chaque enregistrement porte le drapeau \texttt{is\_aggregate} qui en résulte. Les agrégats sont \textbf{conservés dans les séries temporelles} (l'évolution du PIB mondial reste une information utile) mais \textbf{exclus des classements}, du calcul des rangs et des visualisations comparatives, où ils écraseraient les pays. La classification de la Banque Mondiale prévaut en cas de désaccord entre les deux sources, afin d'éviter qu'un territoire dont le code entre en collision avec un groupe FMI (ainsi \texttt{MAF}, Saint-Martin) ne soit écarté à tort.

\subsection{Réconciliation des Libellés et des Codes}
Un même code ISO reçoit un libellé différent selon la source (\textit{Korea, Rep.} à la Banque Mondiale, \textit{Korea} au FMI), et les payloads d'indicateurs du FMI ne transportent pas les libellés. La jointure historique/prévision s'effectue donc sur le \textbf{seul code pays}, le nom de la Banque Mondiale étant retenu comme forme canonique. Sans cette réconciliation, chaque pays se scinde en deux séries — l'une sans prévision, l'autre sans historique — ce qui vide le PIB 2030, le CAGR prévisionnel et le rang projeté.

La jointure sur le code suppose que les deux sources désignent un territoire par le même code. Deux exceptions sont converties vers les codes de la Banque Mondiale: le Kosovo (\texttt{UVK} dans l'API du FMI, \texttt{KOS} dans sa base historique, \texttt{XKX} à la Banque Mondiale) et la Cisjordanie et Gaza (\texttt{WBG}, \texttt{PSE}).

\subsection{Panel de Classement}
Tous les rangs portent sur un même panel: les pays renseignés à la dernière année observée et à l'horizon de prévision, au taux de marché comme à parité (183 pays). Un écart de rang (\texttt{Rank\_Change}, \texttt{Ecart\_Rang\_Nominal\_PPA\_2024}) traduit alors un mouvement, et non l'entrée ou la sortie d'un pays du classement. Classer chaque colonne sur les pays qu'elle couvre faussait 140 écarts de rang sur 183: Taïwan, absent de la Banque Mondiale, entrait 22\textsuperscript{e} au classement 2030; le Pakistan, sans projection du FMI au-delà de 2025, en sortait. Les pays hors panel conservent leurs niveaux de PIB dans la synthèse, sans rang.

\subsection{Portée du Classement à Parité}
Les niveaux de prévision à parité étant obtenus par chaînage à partir du dernier point observé, le classement \texttt{Rank\_PPA} projeté suppose les facteurs de conversion de 2021 inchangés sur tout l'horizon. Or ces facteurs évoluent avec les niveaux de prix relatifs. Le classement projeté à parité doit donc se lire comme l'extrapolation d'une structure de prix figée, et non comme une prévision du pouvoir d'achat relatif à l'horizon: son incertitude excède celle des prévisions de croissance dont il dérive.

\section{Organisation du Dépôt}

\begin{itemize}
    \item \texttt{produire\_rapports.py}, à la racine: point d'entrée unique, qui enchaîne la chaîne complète pour le rapport de référence et, le cas échéant, le plus récent.
    \item \texttt{pib/}: le code, en paquet Python (détail ci-dessous).
    \item \texttt{notebooks/}: les deux notebooks d'analyse (Python et Julia) et leurs générateurs.
    \item \texttt{tests/}: les tests des invariants (pytest).
    \item \texttt{docs/}: la présente documentation, source \LaTeX{} et PDF.
    \item \texttt{data/raw/}: les sources versionnées --- classeur WEO historique, complément lu dans l'API du FMI, et leur notice.
    \item \texttt{data/processed/}, \texttt{data/plus\_recent/}, \texttt{outputs/}: les données et livrables produits, régénérés à chaque exécution et non versionnés.
\end{itemize}

\noindent Modules du paquet \texttt{pib}, chacun exécutable depuis la racine du dépôt par \texttt{python -m pib.<module>}:
\begin{itemize}
    \item \texttt{http\_utils}: requêtes HTTP communes aux collectes (nouvelles tentatives, échec explicite).
    \item \texttt{fetch\_historical\_gdp}: collecte des données historiques (Banque Mondiale).
    \item \texttt{fetch\_forecast\_gdp}: collecte des prévisions (FMI, API DataMapper).
    \item \texttt{gdp\_pipeline}: unification, raccord, volumes, synthèse, choix des rapports et export Excel / CSV.
    \item \texttt{update\_weo\_editions}: éditions récentes du WEO, depuis l'API SDMX du FMI.
    \item \texttt{evaluate\_forecasts}: évaluation des prévisions d'époque du FMI, risque de récession.
    \item \texttt{revisions\_weo}: révisions d'une édition du WEO à la suivante.
    \item \texttt{calibration}: fourchettes et probabilités de récession éprouvées en temps réel.
    \item \texttt{millesimes\_bm}: éditions archivées des WDI de la Banque Mondiale, révisions du réalisé.
    \item \texttt{cas\_de\_crise}: une récession mondiale dans les prévisions du FMI, sur demande (\texttt{-{}-annee 2009}, \texttt{-{}-annee 2020}).
    \item \texttt{visualize\_gdp}: graphiques PNG et tableau de bord HTML.
    \item \texttt{build\_results\_page}: page HTML de résultats, construite depuis les CSV et les PNG produits.
\end{itemize}

\subsection{Bornes Paramétrables}
Les quatre bornes du run (\texttt{-{}-start-year}, \texttt{-{}-end-year}, \texttt{-{}-fcst-start}, \texttt{-{}-fcst-end}) se propagent à l'ensemble de la chaîne: nom du fichier unifié, onglet Excel des séries, années de référence de la synthèse, intitulés de colonnes (\texttt{CAGR\_Historique\_1995\_2023\_Pct} pour un run 1995--2023), frontière historique/prévision et bornes d'axe des graphiques. Ni \texttt{visualize\_gdp.py} ni les notebooks ne codent d'année en dur: les premiers redétectent les bornes dans les données, les seconds reconstruisent les noms de colonnes au chargement.

Sans \texttt{-{}-fcst-end}, l'horizon est celui de l'édition du WEO servie par l'API: la dernière année dont le PIB projeté couvre plus de la moitié des pays (2031 pour l'édition d'avril 2026). Il avance ainsi d'un an à chaque édition de printemps, sans intervention.

La prévision doit commencer l'année suivant la fin de l'historique (\texttt{-{}-fcst-start} vaut par défaut \texttt{-{}-end-year}$+1$): un trou laisserait une année sans donnée, qu'un chaînage franchirait; un chevauchement ferait passer une prévision pour la base observée du raccord. Des bornes incohérentes sont refusées avant toute requête.

Chaque requête est retentée trois fois. Une source ou un indicateur qui reste injoignable interrompt le pipeline avec un code de sortie non nul, sans rien écrire: une colonne manquante viderait silencieusement les métriques qui en dépendent, et une liste de pays manquante laisserait les agrégats entrer dans les classements.

\subsection{Deux Rapports: Référence et Plus Récent}
La Banque Mondiale publie une nouvelle année pays par pays. En septembre 2026, 2025 est disponible pour 186 pays, mais pas pour les Émirats arabes unis (27\textsuperscript{e} économie), les Bahamas ou Aruba: retenir 2025 comme dernière année observée les priverait de rang, s'en tenir à 2024 ignorerait des données publiées.

Sauf si \texttt{-{}-end-year} l'impose, la dernière année observée est donc choisie d'après les données, et deux rapports peuvent coexister:
\begin{itemize}
    \item le \textbf{rapport de référence}, dans \texttt{data/} et \texttt{outputs/}: parmi les cinq dernières années observées, la plus récente dont les pays privés de rang pèsent moins de $0{,}1\%$ du PIB des pays classables sur ces cinq ans --- aujourd'hui 2024, 183 pays classés;
    \item le \textbf{rapport le plus récent}, dans les sous-dossiers \texttt{plus\_recent/}, sur la dernière année publiée --- aujourd'hui 2025, 180 pays classés. Il n'existe que s'il diffère de la référence, et il est supprimé dès que la dernière année publiée devient quasi complète.
\end{itemize}

Le seuil porte sur le PIB plutôt que sur le nombre de pays: 2023 classe un pays de plus que 2024 (Saint-Marin, $0{,}002\%$ du PIB), ce qui ne justifie pas de perdre une année, quand les trois pays absents de 2025 pèsent $0{,}49\%$. Le bloc \texttt{rapport} de chaque \texttt{extraction\_metadata.json} consigne ce choix, et la page de résultats de chaque rapport renvoie vers l'autre. La commande \texttt{produire\_rapports.py} enchaîne collecte, évaluation, graphiques et page de résultats pour chacun.

\section{Évaluation des Prévisions}\label{sec:evaluation}

Confronter une prévision à ce qui s'est produit suppose de disposer de la prévision \emph{telle qu'elle fut publiée}. L'API DataMapper ne le permet pas: elle ne sert que le millésime courant, dont les valeurs passées sont les estimations d'aujourd'hui. Le FMI consolide en revanche l'ensemble de ses éditions depuis 1990 dans la \textit{WEO Historical Forecasts Database} (\texttt{data/raw/WEOhistorical.xlsx}): 73 éditions, 201 entités, années visées de 1988 à 2031.

\subsection{Obtention du Classeur}
Ce fichier n'est pas récupérable par script: le serveur du FMI répond \texttt{403} aux requêtes automatisées, sur toutes les variantes d'URL testées. Il se télécharge depuis un navigateur, à l'adresse \url{https://data.imf.org/en/Search-Results} en recherchant «~WEO Historical Forecasts Database~», puis se dépose dans \texttt{data/raw/} sous son nom d'origine.

Le FMI réorganisant régulièrement ses serveurs et publiant deux éditions par an, tout lien direct est appelé à devenir caduc; la page de recherche reste le point d'entrée stable. Le lien direct en vigueur, la structure détaillée du classeur et la marche à suivre complète figurent dans \texttt{data/raw/LISEZ-MOI-WEOhistorical.md}.

\subsection{Éditions Récentes: l'API SDMX du FMI}
Le téléchargement manuel n'est plus nécessaire pour suivre les nouvelles éditions. L'API SDMX du FMI (\url{https://api.imf.org}), qui accepte les scripts, sert l'édition courante du WEO et quelques éditions archivées, aux valeurs du classeur (vérifié sur les éditions d'avril 2026 et d'octobre 2025, écart inférieur à $0{,}0005$ point) et avec les mêmes codes pays. Le script \texttt{update\_weo\_editions.py} ajoute les éditions absentes du classeur au fichier \texttt{weo\_editions\_api.csv} de \texttt{data/raw/}, cumulatif et versionné que l'évaluation lit en complément; pour une édition présente des deux côtés, le classeur fait foi.

Le flux courant ne précise pas son édition. Elle est déduite de la plus récente archive --- l'édition d'octobre 2025 étant archivée, le flux courant est celle d'avril 2026 --- puis vérifiée: la dernière année servie doit valoir l'année de l'édition plus cinq. Faute de quoi le script s'arrête plutôt que de deviner.

\paragraph{Archive des éditions complètes.} L'API ne sert que l'édition courante et la précédente. Chaque édition servie est donc archivée en entier dans \texttt{data/raw/weo\_archive/} --- tous pays, 145 indicateurs, PIB en dollars courants compris --- telle qu'elle a été servie la première fois, puis jamais réécrite. La base historique ne contenant que des taux, ces archives permettront, édition après édition, de mesurer aussi les erreurs sur le PIB en dollars courants. Les deux premières (octobre 2025 et avril 2026) datent du 29 septembre 2026.

\subsection{Horizon de Projection}
Chaque édition couvre huit années relatives, de $-2$ à $+5$ par rapport à sa date de publication. L'horizon se déduit de l'écart entre l'année visée $y$ et l'année de l'édition $v$:
\begin{equation}
h = y - v
\end{equation}
Un horizon négatif désigne une \emph{ré-estimation} d'une année déjà écoulée, non une prévision: les inclure reviendrait à mesurer la capacité de l'institution à se souvenir du passé. Seules les valeurs telles que $h \geq 0$ sont évaluées.

\subsection{Choix de la Valeur de Référence}
Comparer une projection au chiffre définitif d'aujourd'hui la pénalise pour des révisions statistiques intervenues des années après coup, hors de portée du prévisionniste. La référence retenue est donc la ré-estimation publiée par le FMI un an après l'année visée ($h = -1$, édition d'automne). Pour la croissance du PIB, le calcul est mené en parallèle contre la série observée de la Banque Mondiale, les deux résultats étant conservés: une conclusion qui ne tiendrait qu'à une seule référence ne serait pas robuste. Pour l'inflation et la balance courante, la série du pipeline n'offre pas d'équivalent: seule la référence FMI est calculée.

L'erreur est exprimée en points de croissance et non en écart relatif: rapporter l'erreur d'un taux à sa propre valeur produit des rapports arbitrairement grands dès que ce taux approche de zéro.

La synthèse par horizon donne le biais et l'erreur absolue en moyenne et en médiane. Pour l'inflation, seules les médianes décrivent l'erreur typique: quelques projections d'hyperinflation (le Venezuela à $10\,000\,000\,\%$) portent le biais moyen au-delà de $1\,600$ points à un an, quand le biais médian reste de $-0{,}1$ point. Le script le signale dès que l'erreur absolue moyenne dépasse dix fois la médiane.

\subsection{Résultats: le Biais de Croissance et ses Lectures}
Sur 69\,345 projections de croissance réelle couvrant 196 pays (erreurs en points de croissance):

\begin{center}\small
\begin{tabular}{lrrrcr}
\hline
\textbf{Horizon} & \textbf{Biais moyen} & \textbf{Pondéré PIB} & \textbf{Médiane} & \textbf{IC 95\,\%} & \textbf{Erreur abs.} \\
\hline
0 (année en cours) & $+0{,}18$ & $-0{,}07$ & $0{,}00$ & $[-0{,}03\,;\,+0{,}39]$ & $1{,}70$ \\
1 an & $+1{,}00$ & $+0{,}58$ & $+0{,}34$ & $[+0{,}38\,;\,+1{,}62]$ & $2{,}74$ \\
3 ans & $+1{,}01$ & $+0{,}80$ & $+0{,}60$ & $[+0{,}34\,;\,+1{,}68]$ & $2{,}95$ \\
5 ans & $+0{,}93$ & $+0{,}81$ & $+0{,}62$ & $[+0{,}24\,;\,+1{,}62]$ & $3{,}04$ \\
\hline
\end{tabular}
\end{center}

La prévision de l'année en cours est pratiquement sans biais. Au-delà, la croissance est surestimée, et ce biais ne s'atténue pas avec l'horizon alors que la dispersion, elle, continue de croître. Son ampleur dépend toutefois de la lecture:
\begin{itemize}
    \item le biais moyen compte chaque pays pour un. Pondéré par le PIB de l'année visée --- ce que l'erreur représente pour l'économie mondiale ---, il tombe de $+1{,}0$ à $+0{,}6$ point à un an, et la médiane à $+0{,}3$: quelques fortes surestimations tirent la moyenne;
    \item les récessions mondiales en expliquent une part: sans 2009 ni 2020, le biais moyen à un an serait de $+0{,}62$ point. Rapportée à l'année visée, l'erreur médiane atteint $+7{,}3$ points pour 2020 et $+3{,}9$ pour 2009, et le rebond qui suit est symétriquement sous-estimé;
    \item les pays d'une même année subissent les mêmes chocs, et leurs erreurs ne sont pas indépendantes. L'erreur type du biais moyen est donc calculée par grappes d'années visées, avec la correction $G/(G-1)$; elle est six fois celle qui supposerait les erreurs indépendantes. L'intervalle de confiance à $95\,\%$ exclut zéro sans être étroit;
    \item le biais est plus fort pour les pays pauvres: à un an, $+0{,}7$ point pour les pays à revenu élevé, $+1{,}7$ pour les pays à faible revenu (classification courante de la Banque Mondiale, appliquée à toute la période).
\end{itemize}

\subsection{Avril, Octobre, Prévision Naïve et Croissance Mondiale}
\paragraph{Avril ou octobre.} Pour une même année visée, l'édition d'octobre dispose de six mois d'information de plus que celle d'avril. Sur l'année en cours, elle est sans biais ($+0{,}01$ point, erreur absolue $1{,}40$) quand celle d'avril penche de $+0{,}35$ point (erreur absolue $2{,}01$): l'horizon 0 fait la moyenne de deux situations différentes. L'écart se réduit à un an et disparaît au-delà.

\paragraph{Face à une prévision naïve.} Une prévision utile doit faire mieux qu'une règle simple. La prévision naïve retenue est, pour une édition de l'année $v$, la croissance moyenne des années $v-5$ à $v-2$ telle que ré-estimée par le FMI --- toutes publiées avant l'édition, la ré-estimation de $v-1$ ne paraissant qu'à l'automne de $v$. Le FMI est plus proche du réalisé dans 74\,\% des cas sur l'année en cours (erreur absolue $1{,}59$ contre $3{,}25$), mais dans 56\,\% seulement à cinq ans ($3{,}04$ contre $3{,}70$): à l'horizon des projections, il ne fait que 18\,\% mieux que la règle naïve.

\paragraph{La croissance mondiale.} L'agrégat \og World \fg{} du FMI, écarté de l'évaluation par pays, est évalué à part: biais de $-0{,}10$ point sur l'année en cours, $+0{,}56$ à un an, $+0{,}83$ à cinq ans (intervalle de confiance de $+0{,}20$ à $+1{,}46$). Il l'est contre sa seule ré-estimation: la Banque Mondiale agrège le monde aux taux de change de marché, et non à parité de pouvoir d'achat, et sa croissance mondiale, inférieure de $0{,}4$ point en moyenne (23 années sur 25), gonflerait le biais de ce seul écart de pondération.

\subsection{PIB Prévu et PIB Réalisé}
Une erreur de croissance répétée d'année en année se cumule sur le niveau du PIB. Pour une édition $v$ et un pays, le niveau prévu à l'horizon $h$, rapporté au niveau de l'année précédant l'édition, est le produit des croissances projetées; le niveau réalisé, celui des croissances ré-estimées. Leur rapport donne l'erreur de niveau:
\begin{equation}
e_h = \frac{\prod_{k=0}^{h}\left(1 + g^{\text{prévu}}_{v+k}/100\right)}{\prod_{k=0}^{h}\left(1 + g^{\text{réalisé}}_{v+k}/100\right)} - 1
\end{equation}
Le niveau de départ, commun aux deux, s'élimine. Une année sans croissance réalisée interrompt l'enchaînement.

\begin{center}\small
\begin{tabular}{lrrcrr}
\hline
\textbf{Horizon} & \textbf{Médiane} & \textbf{Pondérée PIB} & \textbf{80\,\% des cas} & \textbf{$>+5\,\%$} & \textbf{$<-5\,\%$} \\
\hline
1 an & $+0{,}4\,\%$ & $+0{,}6\,\%$ & $[-3{,}9\,;\,+6{,}9]$ & $15\,\%$ & $7\,\%$ \\
3 ans & $+2{,}4\,\%$ & $+2{,}4\,\%$ & $[-6{,}2\,;\,+14{,}9]$ & $36\,\%$ & $13\,\%$ \\
5 ans & $+4{,}6\,\%$ & $+4{,}4\,\%$ & $[-7{,}9\,;\,+22{,}7]$ & $49\,\%$ & $15\,\%$ \\
\hline
\end{tabular}
\end{center}

À cinq ans, le niveau prévu dépasse le réalisé de 4 à 5\,\% en médiane, et les erreurs ne se compensent pas: il est trop haut de plus de 5\,\% dans près d'un cas sur deux, trop bas de plus de 5\,\% dans un sur sept. L'erreur médiane va de $+3{,}6\,\%$ pour les pays à revenu élevé à $+6{,}0\,\%$ pour les pays à faible revenu; contre la série de la Banque Mondiale, elle est de $+3{,}6\,\%$. La base historique ne contenant que des taux, cette mesure porte sur le volume: en dollars courants s'ajouteraient les erreurs de change et d'inflation.

\subsection{Efficience: l'Erreur Croît avec la Croissance Annoncée}
Le test d'efficience de Mincer et Zarnowitz ajuste, par horizon, le réalisé sur le prévu: $g^{\text{réalisé}} = a + b\,g^{\text{prévu}}$, sans le 1\,\% de valeurs extrêmes de chaque côté. Une prévision efficace donne $b = 1$. La pente vaut $0{,}93$ sur l'année en cours, $0{,}74$ à un an, $0{,}61$ à trois ans et $0{,}57$ à cinq ans: le réalisé ne suit qu'à moitié les écarts de croissance annoncés d'un pays à l'autre. À cinq ans, le biais médian va de $-0{,}36$ point pour le cinquième des prévisions les plus modestes à $+1{,}35$ pour le cinquième des plus fortes; sur le niveau, l'erreur médiane va de $+1{,}7\,\%$ pour les croissances cumulées projetées les plus faibles à $+6{,}8\,\%$ pour les plus fortes. Les marchés à forte croissance projetée, où se concentre la croissance du trafic aérien, sont ceux dont les projections sont les plus optimistes.

\subsection{Les Projections à l'Aune des Erreurs Passées}
Les projections de l'horizon (2031) viennent de l'édition d'avril 2026, à cinq ans d'horizon. Leur appliquer les erreurs de niveau passées au même horizon donne une fourchette empirique: le niveau réalisé valant le niveau prévu divisé par $(1 + e)$, les 10\textsuperscript{e} et 90\textsuperscript{e} centiles des erreurs passées bornent l'intervalle où seraient tombés 80\,\% des cas comparables. Sont comparables les projections de la même classe de croissance cumulée projetée (cinq classes de même effectif) et du même groupe de revenu; une cellule de moins de cent cas se replie sur la classe seule. Plus la croissance projetée est forte, plus la fourchette penche vers le bas: de $-11{,}9$ à $+6{,}9\,\%$ pour les États-Unis (croissance projetée de $12{,}5\,\%$ sur six ans), de $-21{,}8$ à $+4{,}8\,\%$ pour l'Inde ($46\,\%$).

\paragraph{Calibration.} Calculées sur les éditions 1990-2007, ces fourchettes ont contenu 83\,\% des erreurs des éditions 2008-2019, pour une cible de 80\,\%; elles sont prudentes pour les vingt premières économies (94\,\%) et justes pour les pays à faible revenu (78\,\%). Ce test rétrospectif est refait à chaque exécution. Une première version tirait la fourchette de l'historique propre de chaque pays; elle n'aurait contenu que 67\,\% des erreurs suivantes. Le biais d'un pays ne se reproduit pas d'une période à l'autre: sa corrélation entre 1990-2007 et 2008-2020 est de $0{,}01$.

Ce n'est pas une prévision corrigée, mais la marge d'erreur qu'a connue le FMI; elle porte sur le volume et serait plus large en dollars courants.

\paragraph{En temps réel.} Le module \texttt{calibration} éprouve les fourchettes édition par édition, de 2000 à 2019, chacune calculée sur les seules erreurs connues à sa date (à cinq ans, les éditions jusqu'à $v - 7$). Scores propres (score d'intervalle et CRPS de Gneiting et Raftery, 2007), méthodes comparées sur les mêmes 6\,874 projections, intervalles de confiance par bootstrap en blocs d'années visées. La couverture est de $79\,\%$ (de $76$ à $82\,\%$), de $66$ à $86\,\%$ selon l'édition, et de $58\,\%$ pour des fourchettes tirées de l'historique du pays. Le conditionnement ne se justifie que par le groupe de revenu et pondéré par le PIB: score d'intervalle de $33{,}5$ contre $36{,}0$ pour une fourchette unique, écart significatif; sans pondération, les méthodes poolées se valent.

\subsection{Les Récessions que la Trajectoire ne Montre Pas}
Une trajectoire peut atteindre le niveau prévu en passant par un creux. L'erreur de niveau porte donc aussi, pour chaque édition et horizon $h$, la plus faible croissance projetée et réalisée des horizons 1 à $h$: négative, elle signale au moins une année de recul (croissance annuelle en volume négative, selon la ré-estimation du FMI à un an). L'année de l'édition, en partie connue à sa publication, n'en fait pas partie.

Le FMI n'annonce presque jamais de recul au-delà de l'année en cours: à un an, $2{,}4\,\%$ de ses projections sont négatives, contre $14{,}2\,\%$ des croissances réalisées, et il n'avait annoncé que $10\,\%$ de ces reculs ($59\,\%$ dans l'année même). Sur les cinq années suivant une édition, un recul était annoncé dans $3\,\%$ des cas; il en est survenu un dans $45\,\%$ ($48\,\%$ pondéré par le PIB, $49\,\%$ pour les vingt premières économies, $32\,\%$ hors des périodes contenant 2009 ou 2020), la pire année à $-3{,}4\,\%$ en médiane.

\paragraph{Une probabilité par pays.} Le risque dépend surtout de la croissance projetée: sur cinq ans, un recul est survenu dans $64\,\%$ des cas pour les croissances cumulées projetées les plus faibles, dans $27\,\%$ pour les plus fortes. L'ordre des classes tient avant et après 2007; le niveau dépend des crises de la période ($38\,\%$ des cas pour les éditions jusqu'à 2007, $56\,\%$ ensuite). Calculées sur 1990-2007 et éprouvées sur 2008-2019, les probabilités par classe de croissance et groupe de revenu donnent le meilleur score de Brier ($0{,}250$, contre $0{,}279$ pour une probabilité unique); la fréquence passée des reculs du pays et sa volatilité passée font moins bien qu'une probabilité unique ($0{,}283$). Chaque fourchette reçoit donc, de sa cellule, la probabilité d'au moins une année de recul d'ici son horizon et la profondeur typique de la pire année: $65\,\%$ pour les États-Unis et les grandes économies avancées, $46\,\%$ pour la Chine, $23\,\%$ pour l'Inde, entre 2027 et 2031. Éprouvées en temps réel, ces probabilités ont une compétence de $7{,}2\,\%$ sur une probabilité unique (de $4{,}5$ à $10{,}2\,\%$), due à la classe de croissance; bien ordonnées, elles ont été trop basses sur 2000-2019 ($38\,\%$ pour $49\,\%$ observés).

\subsection{Révisions d'une Édition à l'Autre}
Le module \texttt{revisions\_weo} rapproche chaque prévision de celle de l'édition précédente, pour le même pays et la même année visée, les éditions consécutives seulement (avril puis octobre, octobre puis avril suivant). Les éditions à deux ans et plus de l'année visée ne révisent presque pas (au plus $0{,}03$ point en moyenne); l'édition d'avril de l'année visée retire $0{,}61$ point en moyenne ($0{,}29$ hors 2009 et 2020), celle d'octobre $0{,}32$. Mises bout à bout, les révisions font $-1{,}09$ point, à peu de chose près le biais à cinq ans ($+0{,}93$): il se résorbe dans les dernières éditions.

\paragraph{Test de Nordhaus.} Un prévisionniste qui intègre toute l'information disponible produit des révisions imprévisibles; s'il ne l'intègre que par étapes, une révision en annonce une autre de même sens. La pente de la révision sur la précédente, sans le 1\,\% de valeurs extrêmes de chaque côté, et la corrélation entre les deux restent faibles: de $-0{,}09$ à $+0{,}18$, légèrement positives à l'approche de l'année visée. Ce qui se prévoit, c'est le sens des révisions tardives, pas leur enchaînement.

\paragraph{Ce que change la dernière édition.} L'archive des éditions complètes permet de comparer les deux dernières sur les niveaux. Le raccord, $\text{niveau}_y = \text{observé}_{\text{base}} \times \text{FMI}_y / \text{FMI}_{\text{base}}$, ne retient du FMI que la croissance cumulée projetée: c'est sa révision qui passe dans les projections. Entre octobre 2025 et avril 2026, celle de 2024-2030 change de plus de $2\,\%$ en volume pour 52 pays sur 189, de plus de $5\,\%$ en dollars courants pour 64. Les révisions du niveau de l'historique s'éliminent dans le raccord; elles sont signalées, sur une année observée dans les deux éditions (l'année de l'ancienne édition moins deux): changement d'année de base des prix constants quand le rapport du PIB courant au PIB en volume est révisé de plus de $5\,\%$ (14 pays, dont l'Inde et le Royaume-Uni), nouvelle estimation des comptes quand le PIB en dollars l'est (10 pays). Contrôlé sur une année encore estimée, le test confondait révision de l'inflation et changement de base.

\subsection{Le Réalisé aussi se Révise}
La Banque Mondiale archive chaque édition de ses World Development Indicators (142 éditions depuis 1989, le PIB à partir de 1994). Le module \texttt{millesimes\_bm} les collecte et mesure les révisions depuis la première publication, pour les années 1993 et suivantes: la croissance change de $0{,}13$ point en médiane dans l'année, de $0{,}57$ point à ce jour, et de plus d'un point dans $35\,\%$ des cas; le niveau en dollars est révisé de plus de $10\,\%$ dans $21\,\%$ des cas en cinq ans et $38\,\%$ à ce jour, surtout à la hausse. Contre la croissance que publiait la Banque Mondiale à la fin de l'année suivante, référence indépendante du FMI et connue en temps réel, le biais à un an vaut $+0{,}88$ point, contre $+0{,}86$ face à la ré-estimation du FMI sur les mêmes projections, et $+0{,}68$ face à la série actuelle.

\subsection{Cas d'Étude: la Crise de 2008}
Le module \texttt{cas\_de\_crise} suit une récession mondiale à travers les éditions qui l'encadrent; le détail de 2009 est dans \texttt{docs/cas\_crise\_2008.md}. La chute n'a été vue que dans l'année même: l'édition d'octobre 2008, trois semaines après la faillite de Lehman Brothers, annonçait encore $+3{,}0\,\%$ pour le monde et un recul pour sept pays, quand 89 pays ($77\,\%$ du PIB mondial) ont reculé. Le rebond de 2010 ne l'a pas été davantage: l'édition d'avril 2009 le sous-estime pour $77\,\%$ des pays ($-3{,}0$ points pondéré par le PIB). En niveau, le rebond a relevé la croissance, pas la trajectoire: en 2013, le PIB mondial reste $6{,}4\,\%$ sous le niveau projeté en octobre 2008, celui des économies avancées $7{,}4\,\%$. Le réalisé des niveaux est l'estimation actuelle, tirée de la dernière édition archivée: le classeur ne ré-estime chaque année que deux ans après.

\paragraph{Le trafic aérien.} Avec l'option \texttt{-{}-trafic}, le module confronte le PIB au trafic passagers: Eurostat pour 30 pays européens (trafic par aéroport, homogène depuis 2004), Banque Mondiale pour les États-Unis; la série mondiale de la Banque Mondiale, rompue en 2010 par un changement de couverture, est écartée. En 2013, le trafic accuse un retard de $35\,\%$ sur sa tendance de 2004-2007 en Europe, de $19\,\%$ aux États-Unis; l'écart du PIB à la projection d'octobre 2008, multiplié par une élasticité de $1{,}0$ à $1{,}5$, en explique $30$ à $45\,\%$ des deux côtés. Entre pays européens, la relation entre trafic et PIB, forte avant la crise ($R^2 = 0{,}54$ en 2004-2007), disparaît après ($R^2 = 0{,}06$ en 2007-2013).

\section{Provenance et Reproductibilité}

Chaque exécution écrit \texttt{data/extraction\_metadata.json}: horodatage de l'extraction, bornes du run, indicateurs interrogés, volumes obtenus et couverture par colonne. Ses bornes désignent aussi la série unifiée du dernier run: plusieurs fichiers \texttt{gdp\_unified\_<début>\_<fin>.csv} pouvant coexister, les scripts en aval, les notebooks et les tests lisent celui-là, jamais le plus récemment modifié.

Le champ \texttt{derniere\_mise\_a\_jour} reprend le \texttt{lastupdated} déclaré par l'API de la Banque Mondiale, qui donne le millésime réel des séries historiques --- celles-ci étant révisées, deux extractions à quelques mois d'intervalle ne sont pas interchangeables. L'API DataMapper du FMI n'expose pas le millésime de son \textit{World Economic Outlook}, publié en avril et en octobre: seule la date d'extraction permet de le situer, ce que le fichier indique explicitement plutôt que de laisser la lacune passer pour une absence de révision.

\section{Tests des Invariants}

Le mode de défaillance de ce traitement n'est pas l'exception mais la valeur silencieusement fausse: une colonne vide, un classement trié sur le mauvais indicateur, un taux extrapolé au travers d'une année manquante. La suite \texttt{test\_gdp\_pipeline.py} (173 tests, sans accès réseau) fixe les propriétés correspondantes:

\begin{itemize}
    \item les agrégats sont absents des classements par pays;
    \item un pays ne porte qu'un seul libellé et un seul code sur toute sa série;
    \item la croissance implicite des volumes chaînés égale le taux publié par le FMI;
    \item une année de croissance manquante, ou une année absente, interrompt le chaînage au lieu d'être franchie;
    \item chaque série de niveau est raccordée selon son propre facteur;
    \item tous les rangs portent sur un même panel, et chaque écart de rang égale la différence des rangs affichés;
    \item les intitulés de colonnes suivent les bornes du run, et des bornes incohérentes sont refusées;
    \item chaque rang est cohérent avec le niveau dont il dérive;
    \item un CAGR non calculable vaut \texttt{NaN} et non une valeur de repli;
    \item une source injoignable interrompt la collecte, et une réponse paginée est lue en entier;
    \item la série lue en aval est celle du dernier run;
    \item la référence Banque Mondiale ne sert qu'à évaluer la croissance du PIB;
    \item le rapport de référence reste quasi complet, le plus récent s'y ajoute sans le remplacer, et un rapport devenu sans objet est supprimé;
    \item l'horizon de prévision suit l'édition du WEO, sans prolonger la série par des projections isolées;
    \item une édition du WEO lue dans l'API n'est nommée que si l'horizon servi le confirme, le complément reste cumulatif et le classeur fait foi;
    \item le biais pondéré suit les poids, l'intervalle de confiance est groupé par année visée, et le biais hors récessions les exclut bien;
    \item l'erreur de niveau enchaîne les croissances et s'interrompt au premier réalisé manquant;
    \item la fourchette des projections s'appuie sur la classe de croissance projetée et le groupe de revenu, sa calibration est vérifiée sur une période ultérieure, et ses bornes sont dans le bon sens;
    \item la pente d'efficience et la croissance cumulée projetée sont exactes;
    \item la prévision naïve n'utilise que ce que le FMI savait à l'édition, les éditions d'avril et d'octobre restent distinctes, et la croissance mondiale n'est comparée qu'à la ré-estimation du FMI;
    \item chaque édition archivée l'est une seule fois, telle que servie;
    \item la pire année d'une période exclut l'année de l'édition et s'interrompt au premier réalisé manquant; les premières économies sont comptées par édition;
    \item une révision n'est calculée qu'entre éditions consécutives, et l'historique est contrôlé sur une année observée dans les deux éditions;
    \item le PIB en volume par habitant est le volume divisé par la population raccordée, observé comme projeté, et reste vide faute de population;
    \item une révision du réalisé se mesure depuis la vraie première publication, et la référence en temps réel ne connaît que les éditions parues à la fin de l'année suivante;
    \item le trafic d'Eurostat est attribué au bon pays et à la bonne année, et la part du retard du trafic expliquée par le PIB a le bon signe;
    \item la calibration en temps réel n'utilise aucune erreur encore inconnue à la date de l'édition, et ses scores sont exacts.
\end{itemize}

Les tests marqués \texttt{donnees} portent sur les fichiers réellement produits et se sautent tant que le pipeline n'a pas été exécuté.

La suite a été validée par mutation: chaque défaut correspondant a été réintroduit dans une copie du code afin de vérifier qu'elle échoue effectivement. Tous sont détectés --- un test qui ne tombe jamais ne garantit rien.

\section{Guide des Notebooks}

Deux notebooks, dans le dossier \texttt{notebooks/}, explorent le rapport de référence de manière interactive:
\begin{itemize}
    \item \texttt{gdp\_analysis\_notebook.ipynb}, en Python (pandas, matplotlib, Plotly);
    \item \texttt{gdp\_analysis\_notebook\_julia.ipynb}, en Julia (CSV.jl, DataFrames.jl, Plots.jl).
\end{itemize}
Ils couvrent les mêmes sections:
\begin{enumerate}
    \item \textbf{Chargement \& Inspection} des données unifiées.
    \item \textbf{Visualisation temporelle interactive} avec séparation historique / prévision.
    \item \textbf{Comparaison des CAGR} historique vs prévisionnel (Top 15).
    \item \textbf{Carte de chaleur (Heatmap)} des taux de croissance, des dix dernières années observées à l'horizon.
    \item \textbf{Nominal ou volume}, et classement à parité de pouvoir d'achat.
    \item \textbf{Fonction de recherche interactive} par pays (\texttt{analyze\_country}).
\end{enumerate}

Lancés depuis Jupyter, ils s'exécutent dans \texttt{notebooks/} et retrouvent seuls la racine du dépôt, où sont les données: le notebook Python remonte jusqu'au dossier qui contient le paquet \texttt{pib}, le notebook Julia s'appuie sur l'environnement activé par son noyau (\texttt{-{}-project=@.}), celui du \texttt{Project.toml} de la racine. Leurs titres et libellés suivent les bornes du run. Ils se régénèrent par \texttt{notebooks/create\_notebook.py} et \texttt{notebooks/create\_notebook\_julia.py}.

\end{document}
